\section{The complete coefficient-class theorem}
\label{sec:introduction}

Integer-valued polynomials need not have integer coefficients. In degree
two this distinction changes the rational periodic orbits of conservative
H\'enon maps: the missing coefficient branch admits an exact five-cycle,
whereas monic integral quadratic maps have only periods one through four.
We give the complete rational cycle classification for the whole class
\begin{equation}
 \Int(\Z)_2=\{P\in\Q[t]:\deg P=2,\ P(\Z)\subseteq\Z\},\qquad
 G_P(x,y)=(y,P(y)-x).
 \label{eq:class}
\end{equation}
The domain is all of $\Q^2$. The Jacobian determinant is $+1$,
and a time step always means one application of the displayed map.

For a map $G$ on $\Q^2$ write
\[
 \Per(G,\Q)=\{p\in\Q^2:G^d(p)=p\text{ for some }d\ge1\}.
\]
The least such $d$ is the exact period of $p$.
A coordinate word $(x_0,\ldots,x_{d-1})$ denotes the states
$(x_i,x_{i+1})$, with indices modulo $d$.
We identify words under cyclic rotation only, not under reversal.
Repeated coordinates are allowed. Because adjacent coordinates determine
both past and future, the least period of a valid primitive word is the
least period of its state orbit.

\begin{theorem}[All integer-valued quadratics]
\label{thm:whole_class}
Let $P\in\Int(\Z)_2$. Write uniquely
\begin{equation}
 P(t)=m\frac{t(t-1)}2+nt+r,
 \qquad m,n,r\in\Z,\quad m\ne0.
 \label{eq:newton}
\end{equation}
The following rules give all rational periodic points of $G_P$ and their
exact periods.

If $m=2\kappa$ is even, use the complete monic integral classification
in Theorem~\ref{thm:imported} for
\[
 H(u,v)=(v,v^2+(n-\kappa)v+\kappa r-u),
\]
and divide every coordinate of its cycles by $\kappa$.
If $m$ is odd, put
\begin{equation}
 q=n-\frac{m+1}{2},\qquad
 A=mr+\frac{3q-q^2}{2}.
 \label{eq:odd_parameter}
\end{equation}
Take exactly the cycles of $F_A$ in Theorem~\ref{thm:half_classification}
and replace every coordinate $u$ by $(u-q)/m$.

For every such $P$,
\begin{equation}
 \#\Per(G_P,\Q)\le17.
 \label{eq:bound17}
\end{equation}
Equality holds exactly when $m$ is odd and $A=-1$ in
\eqref{eq:odd_parameter}. The set of possible exact periods, as $P$
ranges over the entire class, is exactly
\begin{equation}
 \{1,2,3,4,5,6,7,9,10\}.
 \label{eq:periodset}
\end{equation}
Original coordinates of periodic points of $G_P$ need not be integers.
\end{theorem}

The new classification needed in this theorem is the half-integral
coefficient family
\begin{equation}
 F_a(x,y)=\left(y,\frac{y(y+1)}2+a-x\right),\qquad a\in\Z.
 \label{eq:F}
\end{equation}
The subscript $a$ in \eqref{eq:F} is an arbitrary integer parameter;
it becomes $A$ when the theorem is applied to \eqref{eq:newton}.

\begin{theorem}[The half-integral normal form]
\label{thm:half_classification}
For every $a\in\Z$, all rational periodic points of $F_a$ are integral.
Its complete oriented cycles are the words in
Tables~\ref{tab:families} and~\ref{tab:exceptions} at parameter $a$.
In Table~\ref{tab:families}, $k$ runs through all nonnegative integers.
At a parameter satisfying multiple rows, take their union.
Every displayed cycle is distinct and has the stated exact period;
there are no omitted degeneracy exclusions.
The sharp bound is $17$, attained exactly at $a=-1$.
\end{theorem}

\begin{table}[htbp]
\centering
\small
\begin{tabular}{@{}L{.27\textwidth}L{.56\textwidth}c@{}}
\toprule
Parameter $a$ & Coordinate words & Period \\
\midrule
$1-k(k+1)/2$ & $(1-k)$, $(k+2)$ & $1$ each \\
$-7-k(k+1)/2$ & $(-k-3,k-2)$ & $2$ \\
$-3-k(k+1)/2$ & $(-k-3,k,k)$, $(k-2,-k-1,-k-1)$ & $3$ each \\
$-1-k(k+1)/2$ & $(-k-1,-k-1,k,k)$ & $4$ \\
\bottomrule
\end{tabular}
\caption{All parametric cycles of $F_a$, with $k\ge0$ in every row.
Words represent their consecutive coordinate pairs and all cyclic
rotations. Different parameter rows may coexist; no cycles are identified
by reversal. Both fixed points and both three-cycles remain distinct
at $k=0$.}
\label{tab:families}
\end{table}

\begin{table}[htbp]
\centering
\small
\begin{tabular}{@{}rll@{}}
\toprule
Parameter $a$ & Exceptional coordinate word & Exact period \\
\midrule
$-12$ & $(-6,-1,-6,4,4)$ & $5$ \\
$-11$ & $(-5,-2,-5,1)$ & $4$ \\
$-8$ & $(-4,-2,-3,-3,-2,-4,0)$ & $7$ \\
$-6$ & $(-3,-2,-2,-3,-1)$ & $5$ \\
$-5$ & $(-4,-1,-1,-4,2,2)$ & $6$ \\
$-3$ & $(-3,-1,0,-2,-2,0,-1,-3,1,1)$ & $10$ \\
$-2$ & $(-2,-1,0,-1,-2,0,0)$ & $7$ \\
$-1$ & $(-2,-1,1,1,-1,-2,1,2,1)$ & $9$ \\
$-1$ & $(-2,0,1,0)$ & $4$ \\
$0$ & $(-1,-1,1,2,2,1)$ & $6$ \\
$0$ & $(-1,0,1,1,0)$ & $5$ \\
\bottomrule
\end{tabular}
\caption{The eleven cycles not in Table~\ref{tab:families}.
Each word is the lexicographically smallest of its rotations.
These are exact finite-certificate outputs, not samples used to infer
the all-parameter theorem.}
\label{tab:exceptions}
\end{table}

\paragraph{What is inherited and what is completed.}
The monic integral classification is the prior internal working
manuscript C412~\cite{c412}; its exact statement is reproduced in
Appendix~\ref{app:integral}. We also use its maximum-coordinate,
annulus, six-symbol, and finite-partial-permutation mechanisms.
The local word argument is reproduced at the point of use, with
attribution. The elementary Newton-basis normalization is not a new
method. The contribution here is the complete missing normal form,
and hence a full natural coefficient-class atlas with a different
sharp bound and period spectrum. The exceptions, equality locus,
and return law are consequences of that one classification.

The all-parameter proof has three parts: no cycles for $a\ge2$ and
two fixed points at $a=1$; a uniform symbolic exhaustion for
$a\le-146$; and exact finite certificates for $-145\le a\le1$,
where $a=1$ is an independently proved control. The finite certificate
is a proof dependency, not an empirical experiment. It was independently
reconstructed using a different graph algorithm before this manuscript
was prepared~\cite{am1certificate}.

\paragraph{Relation to the wider problem.}
Ingram~\cite{ingram2014canonical} develops canonical-height and
specialization results for H\'enon maps and studies a quadratic
rational-period conjecture for $(y,x+y^2+b)$. That displayed map has
Jacobian determinant $-1$, whereas \eqref{eq:class} has determinant
$+1$. Our statement is not a solution of that conjecture or of
uniform boundedness for arbitrary rational-coefficient H\'enon maps.
Kim, Krieger, Postolache, and Szeto~\cite[Theorems A--B]{kim2025many}
construct conservative H\'enon maps with many rational periodic
points and long integer cycles from integer-valued polynomials in
growing odd degree. Their construction motivates keeping the
coefficient class explicit; its degree range is different from the
fixed degree-two classification here. These comparisons assert
scope and dependence, not worldwide priority.
